\documentclass[11pt]{article}

\usepackage[margin=1in]{geometry}
\usepackage{amsmath, amssymb, amsthm}
\usepackage{booktabs}
\usepackage{array}
\usepackage{tabularx}
\usepackage{enumitem}
\usepackage{xcolor}
\usepackage{listings}
\usepackage[hidelinks]{hyperref}
\usepackage{microtype}

\lstset{
  basicstyle=\ttfamily\small,
  breaklines=true,
  columns=fullflexible,
  frame=single,
  framesep=4pt,
  xleftmargin=0pt,
  showstringspaces=false,
}

\newcolumntype{Y}{>{\raggedright\arraybackslash}X}

\title{\textbf{dbemu-bench}: \\
A Benchmark for Database Emulation by Sequence Models}
\author{}
\date{}

\begin{document}

\maketitle

\begin{abstract}
We propose \emph{dbemu-bench}, a benchmark that stress-tests sequence
models and large language models on \emph{emulating a database}: ingest an
initial state, apply a long stream of operations, answer queries. The
benchmark is structured along two orthogonal axes: (i) a
\emph{strategy spectrum} of eleven points ranging from purely
tensor-native execution (parameters, hidden states, KV-cache) through
hybrid prompt-managed memory to agentic delegation to a real database
engine, and (ii) a \emph{regime split} that separates context-limited
transformers (which must adopt explicit memory-management strategies) from
context-unlimited recurrent / state-space architectures (which expose
state-bit-normalized capacity). Tasks are drawn from a 14-family taxonomy
designed so that distinct failure modes (window cliffs, summarization
loss, retrieval miss, planning error) each have at least one task that
isolates them. Metrics are constructed from theoretically grounded
primitives (cell-F1, Earth Mover's Distance, tree edit distance, mutual
information / bits recovered, calibration) so that lossy and lossless
outputs are scored on a common scale and partial correctness is rewarded
on principle rather than by heuristic. The harness exposes a single
\texttt{evaluate} entry point intended to support automated capability
search by AI/ML research agents.
\end{abstract}

\section{Introduction}

A database is, formally, a state machine: a typed store mutated by a
stream of operations and interrogated by queries. ``Emulating'' a
database is therefore a test of long-horizon, structured state tracking
under arbitrary input distributions --- a capability that recurs in
agentic LLM use, in retrieval-augmented systems, and in recent
state-space and recurrent architectures.

This benchmark unifies three concerns:

\begin{enumerate}[leftmargin=*]
  \item \textbf{Architecture}: how much state a model can faithfully
        carry per parameter, per KV byte, and per recurrent state bit.
  \item \textbf{Strategy}: how a context-limited model manages an
        operation stream that exceeds its window --- via eviction,
        summarization, retrieval, or tool delegation.
  \item \textbf{Agency}: how reliably a model orchestrates external
        software (an opaque key-value store, a typed table, a real SQL
        engine) when the database state is externalized.
\end{enumerate}

These three live on the same axis. Section~\ref{sec:spectrum} formalizes
that axis as an eleven-point spectrum; the rest of the paper develops the
inputs (\S\ref{sec:inputs}), operations (\S\ref{sec:ops}), tasks
(\S\ref{sec:tasks}), metrics (\S\ref{sec:metrics}), stress dimensions
(\S\ref{sec:stress}), baselines (\S\ref{sec:baselines}), and a
research-agent API (\S\ref{sec:api}).

\section{Inputs: state representations}
\label{sec:inputs}

The same logical database is rendered four ways so that format
sensitivity is independently measurable:

\begin{description}[leftmargin=*]
  \item[SQLite.] Schema + rows; also the canonical ground-truth executor.
  \item[JSON.] Documents per table, or a single nested object; tests
    nested-state tracking.
  \item[CSV.] Header + rows; tests flat tabular recall.
  \item[FlatBuffers.] Schema-bound binary; tests strict-typed,
    offset-addressed recall.
\end{description}

Each instance is a tuple
$\langle \text{schema}, \text{initial\_state}, \text{op\_stream},
\text{queries}\rangle$. The op stream is format-agnostic at the semantic
level; the harness renders it into each surface form per run.

\section{Operation grammar}
\label{sec:ops}

\begin{description}[leftmargin=*]
  \item[DDL.] \texttt{CREATE / ALTER / DROP TABLE}, \texttt{CREATE INDEX}.
  \item[DML.] \texttt{INSERT}, \texttt{UPDATE}, \texttt{DELETE},
    \texttt{UPSERT}.
  \item[Transaction.] \texttt{BEGIN}, \texttt{SAVEPOINT},
    \texttt{ROLLBACK [TO sp]}, \texttt{COMMIT}.
  \item[Query.] \texttt{SELECT} with filter, join, group-by, aggregate,
    order, limit; \texttt{EXISTS}, \texttt{COUNT}.
  \item[Time-travel.] \texttt{AS OF $\langle$op\_idx$\rangle$} (MVCC
    snapshot reads).
  \item[Bulk.] \texttt{IMPORT csv}, \texttt{MERGE}.
\end{description}

The op stream is the model's input after the initial state; queries are
the evaluation points.

\section{Two model regimes}
\label{sec:regimes}

\paragraph{Regime A --- context-limited transformer.}
Stream length far exceeds the context window. The model must adopt an
explicit memory strategy: sliding KV, summary checkpointing,
hierarchical summarization, retrieval over past operations, external
tool memory, or a learned compressor. Each strategy is a plug-in; the
benchmark scores the \emph{strategy}, not just the underlying model.

\paragraph{Regime B --- context-unlimited / fixed state.}
For RNN, SSM (Mamba), RWKV, or any transformer with full window. Same
tasks, no eviction. Adds a critical axis: state-size-normalized metrics
(accuracy per bit of recurrent state). This is what makes the benchmark
informative for architecture research, not just for prompting research.

\section{The strategy spectrum: full-tensor to agentic-DB}
\label{sec:spectrum}

Eleven points, ordered by how much state leaves the model's tensors and
lives in external software. Each is a valid \texttt{strategy} in the
benchmark; the Pareto sweep across them is the principal research
artifact.

\begin{center}
\small
weights $\rightarrow$ hidden $\rightarrow$ KV $\rightarrow$ scratchpad
$\rightarrow$ opaque KV $\rightarrow$ typed tool $\rightarrow$ SQL
$\rightarrow$ planner+executor $\rightarrow$ pure DB
\end{center}

Along the axis:

\begin{itemize}[leftmargin=*]
  \item \emph{Where state lives}: parameters $\to$ activations $\to$
    KV-cache $\to$ emitted tokens $\to$ external bytes.
  \item \emph{Who interprets it}: the network $\to$ network + prompt
    $\to$ external software.
  \item \emph{What the benchmark measures}: architectural capacity $\to$
    context-management capability $\to$ tool-use / planning capability.
  \item \emph{Failure mode}: forgetting $\to$ eviction lossiness $\to$
    tool-call errors $\to$ planning errors.
\end{itemize}

\subsection{Levels}

\begin{table}[h]
\centering
\small
\begin{tabularx}{\linewidth}{@{}r l l Y@{}}
\toprule
L & Name & State medium & Bench signal isolates \\
\midrule
0  & Parametric            & weights          & training-time memorization     \\
1  & In-context            & KV (full)        & long-context attention         \\
2  & Recurrent             & hidden state     & architecture / state bits      \\
3  & Evicting KV           & bounded KV       & eviction policy                \\
4  & Self-summary          & KV of digests    & model's compression sense      \\
5  & Hierarchical / learned& layered codes    & budgeted compression           \\
6  & Text scratchpad       & emitted tokens   & format discipline              \\
7  & KV tool               & external bytes   & retrieval discipline           \\
8  & Typed tool            & external typed   & schema + planning              \\
9  & SQL tool              & real DB          & SQL competence                 \\
10 & Planner / executor    & DB + plan trace  & meta-reasoning                 \\
11 & Pure DB (oracle)      & DB               & ceiling                        \\
\bottomrule
\end{tabularx}
\caption{The strategy spectrum.}
\end{table}

\paragraph{L0 --- Pure parametric.}
State baked into weights. No runtime update. Tests memorization and
generalization across schemas. Breaks on anything dynamic.

\paragraph{L1 --- Pure in-context.}
Full op-stream re-fed every call. State is the activation pattern induced
by the prompt. Cost is $O(n^2)$ attention; breaks when stream length
exceeds context.

\paragraph{L2 --- Fixed-state recurrent.}
State is the hidden vector. Single pass, no re-read. Tests pure
state-tracking capacity per bit of state --- the architecture signal.
Metrics are reported normalized by state bits.

\paragraph{L3 --- KV-cache with eviction.}
Sliding window, attention-score, or H2O-style policies. Bounded KV
managed by a policy; the locality/recency profile diagnoses the policy.

\paragraph{L4 --- Self-summarization.}
The model emits a natural-language state snapshot at checkpoints; old
tokens dropped, snapshot stays. Tests the model's ability to
\emph{factor} state --- to separate variable from constant, hot from
cold.

\paragraph{L5 --- Hierarchical / learned compressor.}
Multi-level digests (per 10, per 100, per 1000 ops) or a learned encoder
packing rows into a fixed token budget. Tests lossy compression quality
at a fixed budget.

\paragraph{L6 --- Structured in-prompt scratchpad.}
The model maintains an ASCII / Markdown / JSON table in its prompt and
rewrites it each turn. Tests discipline, format adherence, size control.

\paragraph{L7 --- Opaque external KV tool.}
\texttt{kv.put(k,v)}, \texttt{kv.get(k)}, \texttt{kv.scan(prefix)}. No
schema enforcement. Tests when to call, with what keys, and
retrieval recall@k.

\paragraph{L8 --- Typed table tool.}
\texttt{table.create / insert / update / select} over a schema. Engine
enforces types and uniqueness. Tests schema design and query planning,
not state tracking.

\paragraph{L9 --- Full SQL tool.}
The model emits SQL against a real engine (SQLite). The engine
\emph{is} the state. Tests SQL competence, schema understanding,
transactional reasoning.

\paragraph{L10 --- Planner / executor split.}
The model is the planner only; decomposes intent into sub-tool calls
(parser, optimizer, executor, verifier). Tests meta-reasoning and error
recovery.

\paragraph{L11 --- Pure DB passthrough (oracle).}
No model in the loop. Ground-truth upper bound.

\subsection{Why the spectrum matters}

Each level collapses different metrics:

\begin{itemize}[leftmargin=*]
  \item \textbf{L0--L2}: state-bit-normalized accuracy is load-bearing;
    tool-call counts are irrelevant.
  \item \textbf{L3--L6}: locality / recency profile and cost-per-KV-byte
    dominate.
  \item \textbf{L7--L10}: tool-call count, tool-call correctness, plan
    length dominate; raw recall approaches $1.0$ because the database
    does the remembering.
  \item \textbf{L11}: zero variance; pins the ceiling.
\end{itemize}

A complete benchmark run reports a single Pareto frontier across all
eleven levels for each task family. The interesting research questions
live in the gaps:

\begin{itemize}[leftmargin=*]
  \item Where does L2 overtake L3 per byte? The architectural advantage
    of recurrence.
  \item Where does L7 beat L4 per token? Externalization vs.\
    compression.
  \item Where does L9 stop helping? When planning errors eat the
    externalization gains.
\end{itemize}

\section{Task taxonomy}
\label{sec:tasks}

Tasks are grouped by the failure mode they probe.

\begin{table}[h]
\centering
\small
\begin{tabularx}{\linewidth}{@{}l Y Y@{}}
\toprule
ID  & Name                          & Breaks                       \\
\midrule
T1  & Read-after-write fidelity     & weak recall                  \\
T2  & Aggregate fidelity            & lossy summarization noise    \\
T3  & Schema induction              & no schema tracking           \\
T4  & Transactional consistency     & no tentative/committed split \\
T5  & Long-range dependency         & sliding window (cliff)       \\
T6  & Hash-table stress             & summary and RAG              \\
T7  & Counter array                 & both window and summary      \\
T8  & Cross-table join after stream & weak cross-state binding     \\
T9  & Out-of-order replay           & naive autoregressive replay  \\
T10 & MVCC snapshot read            & history discarded            \\
T11 & Type coercion across formats  & non-invariant representations\\
T12 & Adversarial tokenization      & near-duplicate key handling  \\
T13 & Constraint enforcement        & no UNIQUE/FK/CHECK tracking  \\
T14 & Bulk import + diff            & mix of wide-load and update  \\
\bottomrule
\end{tabularx}
\caption{Task families.}
\end{table}

\section{Metrics}
\label{sec:metrics}

All scoring is built from theoretically grounded primitives so that
lossless and lossy outputs are scored on a common scale and partial
correctness is principled.

\subsection{Lossless primitives}
\begin{itemize}[leftmargin=*]
  \item \textbf{Exact match} --- sanity check only.
  \item \textbf{Cell-F1} over $(\text{table}, \text{row\_id},
    \text{col}) \to \text{value}$ tuples, decomposed into row
    precision/recall, column precision/recall, and value accuracy
    conditional on row + column matched.
  \item \textbf{Schema-F1} over $(\text{name}, \text{type})$ pairs.
  \item \textbf{Order metric} (when \texttt{ORDER BY}): Kendall-$\tau$ on
    the returned row sequence vs.\ ground truth.
\end{itemize}

\subsection{Lossy / aggregate primitives}
\begin{itemize}[leftmargin=*]
  \item \textbf{Relative error} for scalar aggregates:
    $|p - t| / \max(|t|, \varepsilon)$.
  \item \textbf{sMAPE} (bounded symmetric variant) for comparison across
    magnitudes.
  \item \textbf{Earth Mover's Distance} for histograms and group-by
    distributions --- a true metric on distributions.
  \item \textbf{Jaccard} and overlap@k for set- and top-k-valued
    answers.
  \item \textbf{Bucketed accuracy} with tolerance bands
    ($1, 5, 10, 25\%$).
\end{itemize}

\subsection{Partial-credit primitives}
\begin{itemize}[leftmargin=*]
  \item \textbf{Tree edit distance} (Zhang--Shasha) over JSON / nested
    rows.
  \item \textbf{Row-tuple edit distance} over tables.
  \item \textbf{Coverage--accuracy decomposition.} Define
    \[
      \mathrm{coverage} \;=\; \frac{|\text{cells\_attempted}|}{|\text{cells\_truth}|}, \quad
      \mathrm{accuracy} \;=\; \frac{|\text{cells\_correct}|}{|\text{cells\_attempted}|}.
    \]
    Report both and their harmonic mean. This prevents
    ``answer-only-what-you-know'' gaming.
\end{itemize}

\subsection{Information-theoretic composite}

Let $T$ be the ground-truth table viewed as a random variable, $P$ the
prediction. Define \textbf{bits recovered}
\[
  I(T; P) \;=\; H(T) \;-\; H(T \mid P),
\]
estimated empirically from the deterministic ground truth. For
deterministic tables this reduces to
\[
  \mathrm{bits\_recovered} \;=\; \sum_{c \in \text{cols}}
    H(c) \cdot \mathrm{cell\_accuracy}(c),
\]
yielding a single scalar that is additive across cells, comparable
across tasks, and naturally normalized by total entropy:
\[
  \tilde{I} \;=\; \frac{I(T;P)}{H(T)} \;\in\; [0, 1].
\]
This $\tilde{I}$ is the default leaderboard scalar.

\paragraph{Calibration.}
When the model emits a per-cell distribution, report Brier score and
expected calibration error (ECE). These separate
\emph{wrong-and-confident} from \emph{wrong-and-uncertain} ---
essential for honest scoring of context-limited regimes, where the model
should know when it has forgotten.

\paragraph{Generative scoring.}
For generative models, report teacher-forced log-likelihood over the
serialized ground truth as an additional, distribution-aware scalar.

\subsection{Cost-aware metrics (for the Pareto)}

\begin{itemize}[leftmargin=*]
  \item bits recovered per prompt token,
  \item bits recovered per KV byte (Regime A) or per state bit (Regime B),
  \item bits recovered per wall-second and per FLOP.
\end{itemize}

\subsection{Locality / recency profiles}

For T5, T6, and T10, plot accuracy vs.\ write-to-query lag. The shape
diagnoses memory strategies:
\begin{itemize}[leftmargin=*]
  \item \emph{Cliff} $\Rightarrow$ sliding window;
  \item \emph{Slow decay} $\Rightarrow$ summarization lossiness;
  \item \emph{Flat} $\Rightarrow$ real memory.
\end{itemize}

\section{Stress dimensions}
\label{sec:stress}

\begin{table}[h]
\centering
\small
\begin{tabularx}{\linewidth}{@{}l Y@{}}
\toprule
Axis & Values \\
\midrule
Stream length         & $10^2, 10^3, 10^4, 10^5, 10^6$            \\
Schema width          & $1, 5, 25, 100$ cols                      \\
Schema depth          & $0, 1, 3, 5$                              \\
Keyspace size         & $10, 10^3, 10^5, 10^7$                    \\
Write churn per key   & $1, 10, 100, 1000$                        \\
Read / write ratio    & $0.01, 0.1, 1, 10$                        \\
Query complexity      & scan, 1-join, 2-join, agg-group, subquery \\
Format                & sqlite, json, csv, flatbuffers            \\
Op ordering           & sorted, random, adversarial-shuffled      \\
\bottomrule
\end{tabularx}
\caption{Stress dimensions swept by the generator.}
\end{table}

A single benchmark run is a seeded sparse-grid sample.

\section{Reference baselines}
\label{sec:baselines}

The harness ships nine reference strategies, one per spectrum region:

\begin{enumerate}[leftmargin=*]
  \item Zero-shot LLM, full stream in context (when it fits).
  \item Sliding-window LLM.
  \item Summary-checkpoint LLM.
  \item RAG-over-ops LLM.
  \item Tool-memory LLM (write/read to a scratch KV store).
  \item Hand-written stateful agent (oracle upper bound for that
    strategy).
  \item RNN baseline (LSTM).
  \item SSM baseline (Mamba).
  \item SQLite ground-truth (perfect-score reference).
\end{enumerate}

\section{Auto-research API}
\label{sec:api}

The harness exposes a single entry point so that an automated research
agent can propose, run, and score a new memory strategy without engaging
with task-specific plumbing:

\begin{lstlisting}
evaluate(strategy_fn, regime, axes_subset) -> {
  scalar,         # normalized bits recovered (default leaderboard)
  breakdown,      # per-task, per-axis dataframe
  pareto_points   # (cost, accuracy) pairs for plotting
}
\end{lstlisting}

The single \texttt{scalar} provides a leaderboard-comparable number;
\texttt{breakdown} provides the diagnostic profile; \texttt{pareto\_points}
support the cost / accuracy frontier that is the principal research
artifact of the benchmark.

\section{Continuous agents and why tensor-native emulation matters}
\label{sec:continuous}

\subsection{Two clocks}

Almost every LLM agent today is \emph{event-driven}: input arrives, a
prompt is built, a forward pass is run, tools are called, a reply is
emitted, the agent goes idle. Between events the model does nothing. It
has no internal clock and cannot consolidate, anticipate,
garbage-collect, or rehearse.

A recurrent or state-space model is structurally different. It runs a
forward pass every tick:
\[
  h_t \;=\; f(h_{t-1}, x_t), \qquad x_t \in \mathcal{X} \cup \{\varnothing\}.
\]
The input $x_t$ is allowed to be empty; the state still advances. The
agent has an internal clock decoupled from input arrival. This is the
substrate on which a \emph{continuous agent} becomes possible:

\begin{itemize}[leftmargin=*]
  \item \emph{Idle compute}: forward passes consume time even when
    nothing arrived; the model can use them to reorganize state.
  \item \emph{Anticipation}: it can predict the next operation and
    pre-compute partial answers before being asked.
  \item \emph{Consolidation}: it can replay and compress recent ops
    into a more compact internal representation --- a sleep / dream
    analog.
  \item \emph{Background maintenance}: it can re-derive invariants,
    refresh derived aggregates, evict stale facts --- exactly the work
    a real database does in the background.
  \item \emph{Adaptive compute}: hard queries get more ticks; trivial
    ones get fewer. Decoupling wall-clock from input is what makes
    adaptive compute naturally expressible.
\end{itemize}

An event-driven transformer can imitate some of this with scratchpads
and tool loops, but only by re-entering the prompt at every step. The
continuous version is cheaper, lower-latency, and qualitatively
different: the agent has \emph{interior time}.

\subsection{A continuous DB-emulating agent}

\begin{lstlisting}
loop t = 0, 1, 2, ...
    x_t  <- next_op_or_none()
    h_t  <- f(h_{t-1}, x_t)
    y_t  <- g(h_t)             # may be empty
    emit(y_t) if salient
end
\end{lstlisting}

Here $f$ is the recurrent update and $g$ a read-out head. The hidden
state $h_t$ \emph{is} the database. When no op arrives, $f$ still runs:
$h_t = f(h_{t-1}, \varnothing)$. That step is not wasted --- it is
where internal reorganization happens.

\begin{table}[h]
\centering
\small
\begin{tabularx}{\linewidth}{@{}Y Y@{}}
\toprule
Idle-tick activity         & Database analog                       \\
\midrule
Re-compress hot keys       & rebuild B-tree pages, vacuum          \\
Refresh derived aggregates & refresh materialized views            \\
Pre-cache likely queries   & query plan cache warm-up              \\
Detect / repair drift      & constraint check, fsck                \\
Rehearse rare facts        & anti-forgetting replay                \\
Predict next op            & speculative execution                 \\
\bottomrule
\end{tabularx}
\caption{Useful uses of no-op forward passes.}
\end{table}

\subsection{Why tensor-native emulation is the load-bearing setting}

The strategy spectrum (\S\ref{sec:spectrum}) runs from L0 (parametric)
through L2 (recurrent / fixed-state) up to L11 (pure external DB). The
agentic levels (L7--L11) are easier and, for engineering, often
preferable. But they are not what a benchmark of \emph{sequence-model
capability} should reward, for three reasons.

\paragraph{Externalization launders the hard part.}
If the model writes \texttt{INSERT INTO users\ldots} to SQLite, the
remembering is done by SQLite. The model only has to parse intent,
generate valid SQL, and read the result back. None of these tests
state-tracking; all of them test translation. A small model with a SQL
tool can score perfectly on a task that a large model without one
cannot touch. The tool, not the model, did the work.

\paragraph{Continuous agency requires internal state.}
A continuous agent runs a forward pass every tick. If its state lives
behind an RPC, every tick that touches state is at least one round-trip.
That cannot be the substrate of an agent ticking per-token or at high
frequency. Continuous agency \emph{requires} state in the tensors.
Tool-memory is fine for episodic recall; it is wrong for moment-to-moment
cognition.

\paragraph{The bench should price what is hard.}
The normalized bits-recovered scalar $\tilde{I}$ is the same number for
L2 and L9, but the cost denominator changes. At L9 an extra bit costs
roughly one extra SQL round-trip; at L2 it costs more recurrent state,
i.e.\ better architecture. The L2 numbers move when architecture moves;
the L9 numbers move when prompting moves. A benchmark that does not
separate these axes will mis-rank progress.

\subsection{The continuous-agent regime}

We add one axis to the stress dimensions of \S\ref{sec:stress}:
\textbf{idle ticks per op}, $\rho \in \{0, 1, 4, 16, 64, 256\}$. For
every op the model receives $\rho$ no-op forward passes before the next
op arrives. Reads are evaluated after the trailing idle window. The
benchmark reports two derivatives:
\begin{itemize}[leftmargin=*]
  \item $\Delta_\rho$ --- accuracy gain per unit idle tick. A purely
    reactive architecture has $\Delta_\rho = 0$; a genuinely continuous
    one has $\Delta_\rho > 0$ and saturates.
  \item Idle-compute efficiency --- bits recovered per idle-tick FLOP.
    Distinguishes architectures that use idle compute well from ones
    that spin.
\end{itemize}

Baselines added for this regime: ACT / PonderNet halting policies;
Mamba / RWKV with looped feedback on $\varnothing$ input; and a
``rumination'' wrapper over an LLM that lets it emit hidden scratchpad
tokens during $\rho$ idle ticks.

\subsection{Why this matters for automated AI research}

If AI research is increasingly to be done by AI research agents, the
benchmark must produce signals those agents can act on: sparse named
failure modes, a single comparable scalar, a Pareto frontier rather
than a leaderboard rank, and a continuous-agent axis that exposes the
least-measured part of the design space: what to do with idle compute.
Tensor-native database emulation, evaluated under both event-driven
and continuous regimes, is the smallest setup that (i) forces all
state into the substrate under study, (ii) exposes architectural
capacity as a measurable quantity, (iii) admits agentic strategies as
one end of a spectrum rather than the only option, and (iv) gives a
place where ``thinking while idle'' is either rewarded or exposed as
decorative.

\section{Four concrete benchmarks}
\label{sec:four}

The design space above is too broad to ship as a single artifact. We
carve it into four concrete benchmarks, each aimed at a distinct
audience but sharing one operation grammar, one oracle (SQLite), and
one JSON wire format $\langle \text{schema}, \text{initial\_state},
\text{ops}, \text{queries}, \text{answers}? \rangle$.

\begin{table}[h]
\centering
\small
\begin{tabularx}{\linewidth}{@{}l Y Y Y@{}}
\toprule
Name & Audience & Shape & Metric \\
\midrule
\texttt{dbemu-mini}   & broad / Kaggle-style        & one dataset, one task, one metric    & cell-F1 \\
\texttt{dbemu-static} & conference, model compare   & frozen subtasks (LongBench / ARC)    & bits recovered + macro \\
\texttt{dbemu-env}    & conference, robustness      & procedural, seeded, hparam grid      & mean $\pm$ std \\
\texttt{dbemu-train}  & from-scratch model research & env + curriculum + reward (SFT + RL) & training library \\
\bottomrule
\end{tabularx}
\caption{Four concrete benchmarks derived from one underlying design.}
\label{tab:four}
\end{table}

\subsection{\texttt{dbemu-mini} --- Kaggle-style}

A drop-in eval anyone can run in five minutes.

\begin{itemize}[leftmargin=*]
  \item \emph{Format.} JSON, one instance per file. Fixed schema:
    three tables, five columns each, INTEGER + TEXT only.
  \item \emph{Initial state.} 100 rows per table.
  \item \emph{Stream.} 1\,000 ops; 50 queries per instance (point-read
    and small range-read only).
  \item \emph{Output.} \texttt{predictions.json}: list of answers in
    query order.
  \item \emph{Metric.} Cell-F1, averaged over queries, averaged over
    instances. One scalar.
  \item \emph{Dataset.} 10\,000 instances; train (5\,k) / public test
    (2.5\,k) / private test (2.5\,k). Frozen forever.
  \item \emph{Env.} Pure offline; one Python scoring script, stdlib
    only.
\end{itemize}

No regime split, no spectrum, no strategy slot. Just one number per
submission. This is the entry-level benchmark; everything else is
optional.

\subsection{\texttt{dbemu-static} --- frozen subtasks}

A citable comparison of off-the-shelf models across distinct capability
slices --- same shape as LongBench or ARC: a frozen test set, a known
taxonomy, a macro-average headline.

\paragraph{Subtasks.}
Five, each isolating one failure mode:
\texttt{T1-read} (read-after-write fidelity),
\texttt{T5-long} (long-range dependency, lag $\geq 10^5$),
\texttt{T6-hash} (random-key hash-table stress),
\texttt{T7-count} (counter-array stress),
\texttt{T10-mvcc} (snapshot read at past op index).

\paragraph{Protocol.}
Splits 800 / 100 / 100 (train / dev / test) per subtask, frozen.
Stream lengths mixed within each subtask: $10^3$, $10^4$, $10^5$.
JSON format only (multi-format is deferred to \texttt{dbemu-env}).
Metric: normalized bits recovered per subtask, macro-average as the
leaderboard scalar.

\paragraph{Two leaderboards.}
\emph{Limited context} (8\,k tokens): models must declare and apply a
named memory strategy. \emph{Unlimited context / recurrent}: no
eviction; results additionally reported per recurrent state bit.

\subsection{\texttt{dbemu-env} --- dynamic procedural env}

Closer to Procgen than to ARC: the dataset is implied by a seeded
generator, not shipped as files. Targets robustness, ablation, and
cliff-finding studies.

\paragraph{Hparam grid.}
Stream length $\in \{10^2, 10^3, 10^4, 10^5, 10^6\}$; schema width
$\in \{1, 5, 25, 100\}$; keyspace $\in \{10, 10^3, 10^5, 10^7\}$;
churn per key $\in \{1, 10, 100, 1000\}$; read/write ratio
$\in \{0.01, 0.1, 1, 10\}$; format
$\in$ \{sqlite, json, csv, flatbuffers\}; idle ticks per op
$\rho \in \{0, 1, 4, 16, 64, 256\}$ (the continuous-agent axis of
\S\ref{sec:continuous}).

\paragraph{Protocol.}
50 seeds per (subtask, hparam point). Report mean $\pm$ std of bits
recovered. Standard outputs: per-axis slice plots, Pareto frontiers
(accuracy vs.\ cost), and the continuous-agent derivative
$\Delta_\rho$.

\subsection{\texttt{dbemu-train} --- training framework (no leaderboard)}

The half of the design space that pure leaderboards cannot measure:
letting a model learn to \emph{be} a database emulator from scratch,
without depending on a pretrained LLM's emergent SQL skill.

\paragraph{Env API (gym-like).}
\begin{lstlisting}
env = DBEmuEnv(seed, hparams)
obs = env.reset()              # initial_state observation
for op in stream:
    obs, info = env.step(op)   # advances oracle, returns next obs
ans = env.query(q)             # returns ground-truth answer
\end{lstlisting}
The env wraps the SQLite oracle in a deterministic stepping interface;
\texttt{obs} is the configured state representation (full table, last
op, idle tick, or empty).

\paragraph{Data generators.}
Infinite procedural stream from the grammar of \S\ref{sec:ops}, with a
curriculum scheduler that stages hparams (small schema $\to$ wide;
short stream $\to$ long; sorted ops $\to$ adversarial).

\paragraph{SFT recipe.}
Targets: per-cell tokens of the post-stream table. Loss:
cross-entropy per cell, optionally curriculum-weighted by difficulty.
Auxiliary: predict-next-op self-supervision.

\paragraph{RL recipe.}
Episode = (stream, queries). Action = emitted answer token sequence.
Reward = bits recovered on the answer, with optional shaping (dense:
per-query cell-accuracy delta; sparse: end-of-episode aggregate).
Off-policy data from SFT generations is supported.

\paragraph{Diagnostics.}
A \emph{state probe} trains a linear head on the hidden state to
predict the ground-truth table; probe accuracy upper-bounds any
read-out head and distinguishes representational-capacity bottlenecks
from read-out bottlenecks. A \emph{capacity curve} sweeps hidden-state
size and reports bits recovered --- the clean architecture signal.

\paragraph{Reference training scripts.}
Small Mamba / RNN / decoder-only transformer trained from random init
on the curriculum. The goal is to establish that the domain is
\emph{learnable from scratch}, not merely amenable to large
pretrained-LLM prompting.

\subsection{Why the four-way split}

\texttt{dbemu-mini} gives a low barrier and broad participation;
\texttt{dbemu-static} provides a stable, citation-friendly comparison;
\texttt{dbemu-env} supports research-grade ablation, robustness, and
the continuous-agent regime; \texttt{dbemu-train} enables architecture
and training-method work independent of frontier-model API access. The
four share one grammar, one oracle, and one wire format, so progress
on any one transfers to the others.

\section{Conclusion}

\emph{dbemu-bench} reduces a heterogeneous research surface --- recurrent
state capacity, KV management, retrieval, tool use, agentic planning ---
to a single benchmark organized around one operation: emulate a database.
The eleven-point strategy spectrum unifies architecture studies and
agent studies under a common scoring function, the information-theoretic
\emph{bits recovered}; the task taxonomy ensures that each canonical
failure mode is isolated by at least one task; and the harness API is
narrow enough to be driven by an automated research agent. The benchmark
should produce, for each model and strategy, a Pareto frontier in the
(cost, fidelity) plane --- and the gaps in that frontier are the open
research questions.

\end{document}
